% ==============================================================================
% CHƯƠNG 3: PHƯƠNG PHÁP NGHIÊN CỨU
% ==============================================================================

\chapter{Phương pháp nghiên cứu}\label{ch:phuong_phap}

\startcontents[chapters]
\printmyminitoc{
Chương này mô tả quy trình nghiên cứu, thiết kế thực nghiệm cùng cách trình bày hình vẽ và bảng số liệu theo chuẩn QĐ 1799.
}

\section{Quy trình triển khai nghiên cứu}\label{sec:ch03_quy_trinh}

Hình vẽ và sơ đồ phải được đặt ngay sau đoạn mô tả (không đặt ngay sau tiêu đề mục), căn giữa và có chú thích cỡ 12bp ở phía dưới (ví dụ Hình \ref{fig:sample_pipeline}).

\begin{figure}[ht!]
  \centering
  \begin{tikzpicture}[>=latex]
    \node[draw, rectangle, rounded corners=3pt, fill=blue!8, minimum width=3.8cm, minimum height=0.9cm] (b1) at (0,0) {1. Thu thập dữ liệu};
    \node[draw, rectangle, rounded corners=3pt, fill=green!10, minimum width=3.8cm, minimum height=0.9cm] (b2) at (4.8,0) {2. Xử lý và Mô hình hóa};
    \node[draw, rectangle, rounded corners=3pt, fill=orange!12, minimum width=3.8cm, minimum height=0.9cm] (b3) at (9.6,0) {3. Đánh giá kết quả};
    \draw[->] (b1) -- (b2);
    \draw[->] (b2) -- (b3);
  \end{tikzpicture}
  \caption{Sơ đồ minh họa quy trình nghiên cứu tổng quát.}\label{fig:sample_pipeline}
\end{figure}

\section{Trình bày Bảng biểu theo quy chuẩn}\label{sec:ch03_bang}

Theo Phụ lục 7 của QĐ 1799, bảng không kẻ đường dọc, chỉ kẻ ngang ở đầu bảng, dưới tiêu đề cột và cuối bảng bằng gói \texttt{booktabs}; tiêu đề bảng đặt phía trên canh trái, cỡ chữ 12bp và nội dung bảng tự động giãn dòng 1.0 (như Bảng \ref{tab:sample_data}).

\begin{table}[ht!]
  \centering
  \caption{Bảng minh họa đặc tả tham số thực nghiệm chuẩn QĐ 1799.}\label{tab:sample_data}
  \begin{tabular}{llr}
    \toprule
    \textbf{Thành phần} & \textbf{Mô tả cấu hình} & \textbf{Giá trị} \\ \midrule
    Tập huấn luyện      & Tỷ lệ phân chia dữ liệu ban đầu & 70,0\% \\
    Tập kiểm thử        & Đánh giá độc lập cuối quy trình & 30,0\% \\
    Số vòng lặp tối đa  & Ngưỡng dừng hội tụ thuật toán & 100 \\
    \bottomrule
  \end{tabular}
\end{table}
